\subsection*{定理 5.8.4.}
\underline{假设 $X$ 满足 (LV) (5.2.2), 且 $\rho$ 满足假设 (A) (5.3). 则对 $\rho$ 的每个非空光滑开集 $U\xrightarrow{\nu}\mathbb P^1$, “模去固定部分”的 Griffiths 同态 (按定义即复合)}
\begin{equation}\tag{5.8.4.1}
\resizebox{\textwidth}{!}{%
\begin{tikzcd}[ampersand replacement=\&,row sep=huge,column sep=huge]
\operatorname{Prim}^n(X,\mathbb Q_\ell(j))
  \arrow[r,"\operatorname{Grif}"]
\&
H^1(U,\nu^*R^{n-1}\rho_*\mathbb Q_\ell(j))
  \arrow[d]
\\
{}
\&
H^1(U,\nu^*E^{n-1}\rho_*\mathbb Q_\ell(j))
\simeq
\dfrac{H^1(U,\nu^*R^{n-1}\rho_*\mathbb Q_\ell(j))}
{(f\times\rho|_U)^*
 H^1(U,H^{n-1}(X,\mathbb Q_\ell(j))_U)}
\end{tikzcd}}
\end{equation}
\underline{为单射.}
\footnote{印本在 (5.8.4.1) 中下方左侧的层上略去了 $\nu^*$, 并未闭合其系数的括号. 公式 (5.8.4.3) 以及 (5.8.4.4) 中同一个下方节点给出改正后的写法.}

\underline{证明.}
由 (5.5.1), 在 $\mathbb P^1$ 上有
\begin{equation}\tag{5.8.4.2}
R^{n-1}\rho_*\mathbb Q_\ell(j)
\simeq
H^{n-1}(X,\mathbb Q_\ell(j))_{\mathbb P^1}
\oplus E^{n-1}\rho_*\mathbb Q_\ell(j),
\end{equation}
而在 $U$ 上有
\begin{equation}\tag{5.8.4.3}
\nu^*R^{n-1}\rho_*\mathbb Q_\ell(j)
\simeq
H^{n-1}(X,\mathbb Q_\ell(j))_U
\oplus \nu^*E^{n-1}\rho_*\mathbb Q_\ell(j).
\end{equation}
这使我们能将 (5.8.4.1) 改写如下:
\begin{equation}\tag{5.8.4.4}
\begin{tikzcd}[row sep=huge,column sep=large]
\operatorname{Prim}^n(X,\mathbb Q_\ell(j))
  \arrow[r,"\operatorname{Grif}"]
&
H^1(U,\nu^*R^{n-1}\rho_*\mathbb Q_\ell(j))
  \arrow[d]
\\
{}
&
H^1(U,\nu^*E^{n-1}\rho_*\mathbb Q_\ell(j)).
\end{tikzcd}
\end{equation}
由 (5.7.1), 映射 (仅由假设 (A) 才得以定义)
\begin{equation}\tag{5.8.4.5}
u:\operatorname{Prim}^n(X,\mathbb Q_\ell(j))
\longrightarrow
H^1(\mathbb P^1,R^{n-1}\rho_*\mathbb Q_\ell(j))
\end{equation}
为单射. 由于 $\mathbb P^1$ 单连通, 分解 (5.8.4.2) 给出
\footnote{印本在 (5.8.4.5) 的靶以及 (5.8.4.6) 的左侧中略去了扭转 $(j)$, 尽管源以及分解 (5.8.4.2) 中保留了它.}
\begin{equation}\tag{5.8.4.6}
H^1(\mathbb P^1,R^{n-1}\rho_*\mathbb Q_\ell(j))
\simeq
H^1(\mathbb P^1,E^{n-1}\rho_*\mathbb Q_\ell(j)).
\end{equation}
因此可以将 (5.8.4.4) 补成交换方块:
\begin{equation}\tag{5.8.4.7}
\begin{tikzcd}[column sep=huge,row sep=huge]
\operatorname{Prim}^n(X,\mathbb Q_\ell(j))
  \arrow[r,"\operatorname{Grif}"]
  \arrow[d,hook,"u"]
&
H^1(U,\nu^*R^{n-1}\rho_*\mathbb Q_\ell(j))
  \arrow[d]
\\
H^1(\mathbb P^1,E^{n-1}\rho_*\mathbb Q_\ell(j))
  \arrow[r,"\nu^*"]
&
H^1(U,\nu^*E^{n-1}\rho_*\mathbb Q_\ell(j)).
\end{tikzcd}
\end{equation}
只须再说明下方箭头为单射. 这由下式
\begin{equation}\tag{5.8.4.8}
E^{n-1}\rho_*\mathbb Q_\ell(j)
\simeq
\nu_*\nu^*E^{n-1}\rho_*\mathbb Q_\ell(j)
\end{equation}
及下面的熟知引理得出.

\subsection*{引理 5.8.5.}
\underline{设 $\nu:A\to B$ 为拓扑斯的态射, $F$ 为 $A$ 上的层. 则典范映射}
\begin{equation}\tag{5.8.5.1}
H^1(B,\nu_*F)\longrightarrow H^1(A,F)
\end{equation}
\underline{为单射.}

\underline{证明.}
对第一象限中的谱序列,
\begin{equation}\tag{5.8.5.2}
E_2^{1,0}\simeq E_\infty^{1,0}.
\end{equation}
因而, 对 Leray 谱序列
\[
E_2^{p,q}=H^p(B,R^q\nu_*F)
\Longrightarrow H^{p+q}(A,F),
\]
有
\begin{equation}\tag{5.8.5.3}
E_2^{1,0}=H^1(B,\nu_*F)
\simeq E_\infty^{1,0}
=F^1H^1(A,F).
\end{equation}
取 $F=\nu^*E$, 其中 $E$ 为 $B$ 上的层, 得
\begin{equation}\tag{5.8.5.4}
H^1(B,\nu_*\nu^*E)
\lhook\joinrel\longrightarrow
H^1(A,\nu^*E).
\end{equation}
